\chapter{Présentation du projet}

Ce chapitre situe le besoin, les utilisateurs visés, les contraintes du domaine automobile et la démarche qui a conduit à la solution finale.


\section{Contexte}

En Tunisie, les échanges de pièces d'occasion passent souvent par des annonces généralistes, des groupes sociaux ou des contacts directs. Ces canaux favorisent la proximité, mais ils structurent peu les informations essentielles. Le titre d'une annonce peut mentionner un véhicule sans préciser la génération, le moteur ou la position. Les coordonnées du vendeur sont parfois exposées publiquement, tandis que les mécanismes de réputation et de modération restent limités.


E-biessa centralise ces échanges dans une application spécialisée. Le catalogue impose des champs cohérents, permet plusieurs véhicules compatibles et distingue la référence OEM principale des références interchangeables. Les fonctionnalités de confiance - validation des vendeurs, avis, signalements, suspension et journal d'audit - accompagnent la recherche et la transaction.


\section{Étude de l'existant}

{\small\singlespacing\arrayrulecolor{tableline}
\begin{longtable}{|>{\raggedright\arraybackslash}p{\dimexpr 0.240\textwidth-2\tabcolsep-2\arrayrulewidth\relax}|>{\raggedright\arraybackslash}p{\dimexpr 0.320\textwidth-2\tabcolsep-2\arrayrulewidth\relax}|>{\raggedright\arraybackslash}p{\dimexpr 0.440\textwidth-2\tabcolsep-2\arrayrulewidth\relax}|}
\hline
\rowcolor{tablehead} \color{white}\bfseries Approche & \color{white}\bfseries Forces & \color{white}\bfseries Limites observées \\ \hline
\endfirsthead
\rowcolor{tablehead} \color{white}\bfseries Approche & \color{white}\bfseries Forces & \color{white}\bfseries Limites observées \\ \hline
\endhead
\hline\endfoot
Petites annonces généralistes & Audience large, publication rapide & Compatibilité peu structurée, filtres insuffisants \\ \hline
\rowcolor{tablerow} Groupes sociaux & Échanges directs, communauté locale & Recherche instable, données dispersées, modération variable \\ \hline
Catalogues professionnels & Références précises, stocks organisés & Souvent centrés sur le neuf et sur un inventaire fermé \\ \hline
\rowcolor{tablerow} E-biessa & Domaine spécialisé, recherche hybride, confiance intégrée & Dépend de la qualité des données saisies et validées \\ \hline
\end{longtable}
}
% End of chapter.
\section{Problématique}

\textbf{Question directrice}

Comment offrir une recherche naturelle et rapide de pièces d'occasion, tout en maintenant une compatibilité vérifiable, une identité fiable des utilisateurs et des outils de modération adaptés ?


Trois difficultés structurent le projet. Premièrement, les données véhicule doivent être propres et reliées à chaque annonce sans mélanger des caractéristiques provenant de véhicules différents. Deuxièmement, les rôles doivent limiter chaque action : un acheteur ne publie pas avant validation comme vendeur, un vendeur ne modère pas, et un administrateur dispose d'un historique d'actions. Troisièmement, l'intelligence artificielle doit améliorer la compréhension sans remplacer les règles métier.


\section{Solution proposée}

\begin{itemize}
\item une application Angular responsive, utilisable en thème clair ou sombre sans modification de géométrie ;

\item un noyau Spring Boot sécurisé pour les comptes, annonces, images, messages, avis et opérations administratives ;

\item un service FastAPI pour l'indexation, la recherche vectorielle et l'extraction de contraintes ;

\item une base MySQL comme source de vérité et un index Qdrant reconstruit à partir des annonces actives ;

\item un modèle Qwen adapté au vocabulaire français, anglais et tunisien, encadré par des validations déterministes ;

\item un moteur de recommandations explicable fondé sur les préférences véhicule et les interactions privées ;

\item une couche agentique en lecture seule pour orchestrer recherche, compatibilité et recommandation.
\end{itemize}

\section{Démarche de réalisation}

La réalisation a suivi une démarche incrémentale. Chaque fonctionnalité a été développée, vérifiée et documentée avant l'introduction des fonctions dépendantes. Le socle marketplace a ainsi précédé la recherche sémantique, la personnalisation et les aides multimodales. Cette organisation ne suppose pas la tenue de cérémonies Scrum formelles.


{\small\singlespacing\arrayrulecolor{tableline}
\begin{longtable}{|>{\raggedright\arraybackslash}p{\dimexpr 0.240\textwidth-2\tabcolsep-2\arrayrulewidth\relax}|>{\raggedright\arraybackslash}p{\dimexpr 0.320\textwidth-2\tabcolsep-2\arrayrulewidth\relax}|>{\raggedright\arraybackslash}p{\dimexpr 0.440\textwidth-2\tabcolsep-2\arrayrulewidth\relax}|}
\hline
\rowcolor{tablehead} \color{white}\bfseries Incrément & \color{white}\bfseries Contenu principal & \color{white}\bfseries Critère de sortie \\ \hline
\endfirsthead
\rowcolor{tablehead} \color{white}\bfseries Incrément & \color{white}\bfseries Contenu principal & \color{white}\bfseries Critère de sortie \\ \hline
\endhead
\hline\endfoot
Socle & Authentification, rôles, profils, catalogue et annonces & Parcours acheteur/vendeur cohérents \\ \hline
\rowcolor{tablerow} Confiance & Validation vendeur, avis, signalements, administration & Actions auditées et autorisées \\ \hline
Communication & Contact protégé, messagerie et images & Échange temps réel entre deux comptes \\ \hline
\rowcolor{tablerow} IA & Embeddings, Qdrant, Qwen adapté et recherche ancrée & Résultats compatibles et explicables \\ \hline
Personnalisation & Préférences, historique minimal et agents & Suggestions privées, contrôlables et justifiées \\ \hline
\rowcolor{tablerow} Validation & Tests unitaires, intégration et E2E & Scénarios critiques reproductibles \\ \hline
\end{longtable}
}

